\section{Introduction}
\label{sec:intro}

The second-hand car market is one of the largest consumer goods markets globally, with millions of transactions per year.
Accurately estimating a fair market price benefits both buyers—who risk overpaying—and sellers—who risk underpricing their vehicles.
Traditional valuation relies on domain expertise and lookup tables, but these approaches fail to capture the complex, non-linear interactions among vehicle attributes such as brand, age, mileage, fuel type, and transmission.
Machine learning offers a data-driven alternative that can model these interactions directly from historical listing data \cite{gegic2019car}.

This project investigates machine learning methods for used car price prediction using the UK Used Car Dataset \cite{usedcardataset2021}, a publicly available collection of $\sim$108{,}000 listings scraped from a major UK automotive marketplace.
The dataset provides a rich combination of numerical and categorical features, real-world missing values, and a heavily right-skewed continuous target (price in GBP), making it well-suited for exploring both modelling and preprocessing questions.

\paragraph{Research Questions.}
We address four specific research questions (RQs):
\begin{enumerate}
  \item[\textbf{RQ1}] To what extent can machine learning models accurately predict used car prices from available features?
  \item[\textbf{RQ2}] How do models spanning the bias–variance spectrum (linear models, instance-based, tree-based, and neural networks) compare in predictive performance and generalisation?
  \item[\textbf{RQ3}] How do different data preprocessing strategies—target transformation, missing value imputation, and feature scaling—affect model performance?
  \item[\textbf{RQ4}] Which features contribute most significantly to price prediction, and how can their importance be interpreted?
\end{enumerate}

RQ1 establishes a baseline understanding of task difficulty. RQ2 provides a systematic model comparison linked to the bias–variance trade-off. RQ3 directly examines the downstream impact of preprocessing decisions, with particular relevance to distance-sensitive models such as KNN. RQ4 connects model behaviour to domain knowledge about the used car market.

The remainder of the paper is structured as follows: Section~\ref{sec:literature} reviews related work; Section~\ref{sec:method} describes our dataset, preprocessing pipeline, models, and evaluation setup; Section~\ref{sec:results} presents experimental results; Section~\ref{sec:discussion} analyses findings; and Section~\ref{sec:conclusion} concludes.
